\id{IRSTI 27.47.23, 28.23.29, 15.21.69}{}

\begin{header}
\swa{A.T.~Mazakova, Sh.A.~Jomartova, T.Zh.~Mazakov, B.~Zhao, A.M.~Abubakirov}{ARCHITECTURE AND APPLICATION OF A BIOTECHNICAL SYSTEM FOR HUMAN PSYCHOPHYSIOLOGICAL TESTING}

\tsp{1}A.T.~Mazakova\alink{https://orcid.org/0000-0003-3019-3352},
\tsp{1}Sh.A.~Jomartova\alink{https://orcid.org/0000-0002-5882-5588},
\tsp{1,2}T.Zh.~Mazakov\alink{https://orcid.org/orcid=0000-0001-9345-5167}\envelope,
\tsp{3}B.~Zhao\alink{https://orcid.org/0009-0002-5572-7306},
\tsp{1}A.M.~Abubakirov\alink{https://orcid.org/0009-0007-3664-9185}
\end{header}

\begin{affil}
\tsp{1}Al-Farabi Kazakh National University, Almaty, Kazakhstan,

\tsp{2}International University of Engineering and Technology, Almaty, Kazakhstan,

\tsp{3} Dalian Polytechnic University, Dalian, China

\corrauthor{Corresponding author: tmazakov@mail.ru}
\end{affil}

This article examines the architecture and practical application of a
biotechnical system for human psychophysiological testing, designed to
objectively assess the functional and psychoemotional state of an
individual during psychological testing. The relevance of this research
stems from the increasing role of the human factor in ensuring the
reliability and safety of modern technical and organizational systems,
as well as the need to integrate instrumental diagnostic methods into
professional selection, personnel certification, and functional state
monitoring.

The proposed biotechnical system implements an integrated approach
combining hardware recording of physiological signals
(electrocardiograms and photoplethysmograms) and software automation of
psycho\-logical testing, followed by intelligent data processing. The
system' s architecture is modular and includes a hardware
subsystem for recording biomedical signals, a microprocessor-based
primary processing mo\-dule, a data transmission interface, and a software
platform for analysis, visualization, and reporting.

The study developed algorithms for digital processing of biomedical
signals, including multi-stage filtering, artifact suppression, cardiac
cycle segmentation, detection of characteristic ECG points,
photop\-lethysmogram analysis, calculation of heart rate variability
parameters, wavelet analysis of non-stationary components, and the
generation of an extended vector of diagnostic features. An ensemble
classification model is used for intelligent data interpretation,
ensuring robust recognition of psychophysiological states in the face of
interindividual variability in physiological parameters.

The results of the experimental testing confirm the viability of the
proposed architecture and the practical applicability of the
biotechnical system for tasks of professional selection, personnel
certification, and functional monitoring in technically complex and
socially responsible areas of activity.

{\bfseries Keywords\emph{: }}biotechnical system, psychophysiological
testing, electrocardiogram, photoplethysmogram, heart rate variability,
digital signal processing, intelligent diagnostics.

\begin{header}
АДАМНЫҢ ПСИХОФИЗИОЛОГИЯЛЫҚ ТЕСТІЛЕУІНЕ АРНАЛҒАН БИОТЕХНИКАЛЫҚ ЖҮЙЕНІҢ АРХИТЕКТУРАСЫ ЖӘНЕ ҚОЛДАНЫЛУЫ

\tsp{1}А.Т.~Мазакова,
\tsp{1}Ш.А.~Джомартова,
\tsp{1,2}Т.Ж.~Мазаков\envelope,
\tsp{3}Б.~Чжао,
\tsp{1}А.М.~Абубакиров
\end{header}

\begin{affil}
\tsp{1}Әл-Фараби атындағы Қазақ Ұлттық Университеті, Алматы, Қазақстан,

\tsp{2}Халықаралық инженерлік-технологиялық университет, Алматы, Қазақстан,

\tsp{3}Далянь политехникалық университеті, Далянь, Қытай,

e-mail: tmazakov@mail.ru
\end{affil}

Бұл мақалада психологиялық тестілеу кезінде жеке тұлғаның функционалдық
және психоэмоционалдық жағдайын объективті бағалауға арналған адамның
психофизиологиялық тестілеуіне арналған биотехникалық жүйенің
архитектурасы мен практикалық қолданылуы қарастырылады. Бұл зерттеудің
өзектілігі заманауи техникалық және ұйымдастырушылық жүйелердің
сенімділігі мен қауіпсіздігін қамтамасыз етудегі адам факторының рөлінің
артуынан, сондай-ақ аспаптық диагностикалық әдістерді кәсіби іріктеуге,
персоналды сертификаттауға және функционалдық жағдайды бақылауға
біріктіру қажеттілігінен туындайды.

Ұсынылған биотехникалық жүйе физиологиялық сигналдарды
(электрокардиограммалар мен фотоплетизмограммалар) аппараттық жазуды
және психологиялық тестілеуді бағдарламалық автоматтандыруды
біріктіретін интеграцияланған тәсілді жүзеге асырады, содан кейін
интеллектуалды деректерді өңдеу жүзеге асырылады. Жүйенің архитектурасы
модульдік болып табылады және биомедициналық сигналдарды жазуға арналған
аппараттық ішкі жүйені, микропроцессорға негізделген бастапқы өңдеу
модулін, деректерді беру интерфейсін және талдау, визуализациялау және
есеп беру үшін бағдарламалық платформаны қамтиды. Зерттеу барысында
биомедициналық сигналдарды сандық өңдеу алгоритмдері әзірленді, оның
ішінде көп сатылы сүзгілеу, артефактіні басу, жүрек циклін сегменттеу,
ЭКГ-ға тән нүктелерді анықтау, фотоплетизмограмманы талдау, жүрек соғу
жиілігінің өзгергіштік параметрлерін есептеу, стационар емес
компоненттердің толқындық талдауы және диагностикалық белгілердің
кеңейтілген векторын жасау. Ансамбльдік жіктеу моделі физиологиялық
параметрлердегі жеке аралық өзгергіштік болған кезде психофизиологиялық
күйлерді сенімді тануды қамтамасыз ететін интеллектуалды деректерді
түсіндіру үшін қолданылады.

Эксперименттік сынақ нәтижелері ұсынылған архитектураның өміршеңдігін
және биотехникалық жүйенің техникалық тұрғыдан күрделі және әлеуметтік
жауапты салаларда кәсіби іріктеу, персоналды сертификаттау және
функционалдық мониторинг үшін практикалық қолданылуын растайды.

{\bfseries Түйін сөздер:} биотехникалық жүйе, психофизиологиялық тестілеу,
электрокардиограмма, фотоплетизмограмма, жүрек соғу жиілігінің
өзгергіштігі, сандық сигналдарды өңдеу, интеллектуалды диагностика.

\begin{header}
АРХИТЕКТУРА И ПРИМЕНЕНИЕ БИОТЕХНИЧЕСКОЙ СИСТЕМЫ ПСИХОФИЗИОЛОГИЧЕСКОГО ТЕСТИРОВАНИЯ ЧЕЛОВЕКА

\tsp{1}А.Т.~Мазакова,
\tsp{1}Ш.А.~Джомартова,
\tsp{1,2}Т.Ж.~Мазаков\envelope,
\tsp{3}Б.~Чжао,
\tsp{1}А.М.~Абубакиров
\end{header}

\begin{affil}
\tsp{1}Казахский национальный университет им. аль-Фараби, Алматы, Казахстан,

\tsp{2}Международный инженерно-технологический университет, Алматы, Казахстан,

\tsp{3}Даляньский политехнический университет, Далянь, Китай,

e-mail: tmazakov@mail.ru
\end{affil}

В статье рассматривается архитектура и практическое применение
биотехнической системы психофизиологического тестирования человека,
предназначенной для объективной оценки функционального и
психоэмоционального состояния личности в процессе выполнения
психологических тестовых заданий. Актуальность исследования обусловлена
возрастающей ролью человеческого фактора в обеспечении надёжности и
безопасности современных технических и организационных систем, а также
необходимостью внедрения инструментальных методов диагностики в процессы
профессионального отбора, аттестации персонала и мониторинга
функционального состояния человека.

Предложенная биотехническая система реализует интегрированный подход,
объединяющий аппаратную регистрацию физиологических сигналов
(электрокардиограммы и фотоплетизмограммы) и программную автоматизацию
психологического тестирования с последующей интеллектуальной обработкой
данных. Архитектура системы построена по модульному принципу и включает
аппаратную подсистему регистрации биомедицинских сигналов,
микропроцессорный модуль первичной обработки, интерфейс передачи данных
и программную платформу анализа, визуализации и формирования
аналитических отчётов.

В работе разработаны алгоритмы цифровой обработки биомедицинских
сигналов, включающие многоэтапную фильтрацию, подавление артефактов,
сегментацию кардиоциклов, детектирование характерных точек ЭКГ, анализ
фотоплетизмограммы, вычисление параметров вариабельности сердечного
ритма, вейвлет-анализ нестационарных компонент и формирование
расширенного вектора диагностических признаков. Для интеллектуальной
интерпретации данных применяется ансамблевая модель классификации,
обеспечивающая устойчивое распознавание психофизиологических состояний в
условиях межиндивидуальной вариабельности физиологических параметров.

Результаты экспериментальной апробации подтверждают работоспособность
предложенной архитектуры и практическую применимость биотехнической
системы для задач профессионального отбора, аттестации персонала и
функционального мониторинга в технически сложных и социально
ответственных областях деятельности.

{\bfseries Ключевые слова:} биотехническая система, психофизиологическое
тестирование, электрокардиограмма, фотоплетизмограмма, вариабельность
сердечного ритма, цифровая обработка сигналов, интеллектуальная
диагностика.

\begin{multicols}{2}
{\bfseries Introduction.} The current stage of scientific and technological
progress is characterized by rapidly increasing levels of automation,
increasing complexity of technical systems, and increasing volumes of
processed information. Under these conditions, the role of the human
factor in ensuring the reliability, safety, and efficiency of complex
ergatic and cyber-physical systems is increasing. Despite the widespread
implementation of intelligent control algorithms, final decision-making
in many critical areas still rests with the human operator,
necessitating constant monitoring of their functional and
psycho-emotional state.

In aviation, transport, energy, telecommunications, industry, law
enforcement and the military, the professional activities of personnel
are associated with high responsibility, the need to quickly respond to
emergency situations and process large volumes of information under time
pressure {[}1,2{]}. In such conditions, decreased concentration,
fatigue, emotional stress and disruption of psychophysiological
adaptation mechanisms can lead to operator errors and, as a consequence,
to emergency situations with severe socio-economic consequences.

In this regard, the task of objectively assessing a
person' s functional state and their ability to
effectively perform professional duties under conditions of increased
psycho-emotional stress is particularly relevant {[}3{]}. Traditional
methods of psychological diagnostics based on questionnaires, surveys,
and expert assessments have a number of limitations due to the
subjectivity of self-reports, the influence of the
subject' s motivational attitudes, and social
desirability effects. Furthermore, such methods do not allow for the
detection of latent stress reactions manifesting at the level of the
autonomic nervous system and cardiovascular regulation.

One promising area is the development of biotechnical systems that
integrate a biological object and a technical device into a single
system for recording, processing, and interpreting information. Such
systems allow for the integration of psychological methods with
hardware-based recording of physiological parameters, ensuring greater
diagnostic reliability.

In recent years, instrumental methods of psychophysiological diagnostics
based on the recording of objective physiological indicators reflecting
the state of the body' s regulatory systems have become
increasingly widespread {[}4{]}. These indicators include
electrocardiograms, photoplethysmograms, galvanic skin responses,
respiratory parameters, electroencephalograms, and other biosignals.
Their analysis allows us to identify signs of sympathoadrenal system
activation, assess autonomic balance, and diagnose the degree of
psychoemotional stress.

The electrocardiogram is of particular interest as one of the most
informative and accessible biomedical signals {[}5-7{]}. Heart rate and
heart rate variability parameters are widely used in medical and
psychophysiological practice to assess the body' s
functional state, stress resistance, and adaptive reserves. Analysis of
the temporal and spectral characteristics of the ECG allows for the
acquisition of objective information on the state of central and
autonomic regulatory mechanisms.

At the same time, recording physiological parameters in isolation
without considering the psychological context does not allow for a full
interpretation of the observed responses. To obtain a reliable
psychophysiological picture, it is necessary to correlate the
body' s physiological responses with specific cognitive
and emotional stimuli. In this regard, integrating psychological testing
with the synchronous recording of physiological signals is a promising
approach.

This integrated approach is implemented within biotechnical systems,
which are complex ergatic assemblies that unite a biological object (a
human) and technical means for recording, processing, and interpreting
information. A biotechnical psychophysiological testing system forms a
closed loop: "human - measurement system - analytical system --
specialist," ensuring objective and reproducible diagnostics of
functional state.

Currently, various hardware and software systems for psychophysiological
diagnostics exist worldwide, including polygraph systems and functional
monitoring systems {[}8,9{]}. However, most of them are focused either
on lie detection or medical diagnostics and do not provide close
integration with psychological testing for professional selection and
personnel certification. Furthermore, many foreign solutions are
expensive, require specialized maintenance, and are not always adapted
to the operating conditions of departmental and corporate information
systems.

In this regard, the development of domestic biotechnical
psychophysiological testing systems focused on professional selection,
pre-shift and pre-trip inspection, functional monitoring, and scientific
research is urgent. Such systems must combine reliable hardware,
flexible software, high diagnostic information yield, and ease of use.

This paper is devoted to the development of an architecture and analysis
of the practical application of a biotechnical system for human
psychophysiological testing, implementing an event-driven approach to
the integration of psychological testing and segmented electrocardiogram
recording {[}10{]}. The paper examines the principles of system
construction, the organization of information flows, algorithms for
digital processing of biomedical signals, and methods for intelligent
data interpretation.

The aim of the study is to create a universal biotechnical platform that
provides an objective assessment of the functional and psycho-emotional
state of a person during the performance of psychological test tasks and
is suitable for widespread implementation in applied and departmental
information environments.

{\bfseries Materials and methods.} The biotechnical system of
psychophysiological testing being developed is intended for the
automated conduct of psychological tests with the simultaneous recording
of the physiological reactions of the subject' s body and
belongs to the class of ergatic biotechnical systems and is intended to
automate the procedure of psychophysiological testing with the
synchronous recording of the physiological reactions of the body.

Three people participate in the testing process: 1) the test subject, 2)
the BTS operator, and 3) the psychologist. Figures 1-2 show the
interaction diagrams.

\fig[0.8\columnwidth]{i/image14}[Fig.1- Scheme of interaction between the BTS operator and the
subject]

\fig[0.8\columnwidth]{i/image15}[Fig.2 - Scheme of processing testing results by a psychologist]

The system architecture is built on a multi-level modular principle and
includes two key functional blocks:

1) a psychological testing block;

2) a block for recording, synchronizing, and writing the
electrocardiogram to a results file.

In general, the system consists of the following functional levels: 1)
recording of biomedical signals, 2) analog-to-digital conversion, 3)
microprocessor processing, 4) data transmission and storage, and 5)
intelligent analysis and visualization.

The system includes:

1) ECG and photoplethysmogram sensors;

2) analog-to-digital converter;

3) microcontroller module;

4) USB interface;

5) personal computer with analysis software platform;

6) results database and report generation module.

The integration of these blocks ensures the formation of an objective
psychophysiological profile of the individual and allows for the
analysis of physiological reactions to each individual stimulus (test
question).

{\bfseries Results and discussion.} The psychological testing unit
automates the psychodiagnostic examination procedure and ensures
standardized presentation of stimulus material, recording of the
subject' s responses, and recording of the time
parameters for completing tasks.

\emph{The unit performs the following main functions:}

- formation of a battery of psychological tests;

- sequential presentation of questions (stimuli);

- recording the start and end of each question;

- recording the selected answer;

- measuring the reaction time of the subject;

- synchronization with the ECG recording unit.

\emph{The testing process includes the following steps:}

- loading the test battery from the database;

- initializing the testing session;

- displaying instructions to the subject;

- sequential presentation of questions;

- recording responses and reaction times;

- transferring time stamps to the ECG recording unit;

- generating a psychological testing protocol.

Each test question is considered as a separate psychophysiological
stimulus, to which the subject's body forms an individual physiological
response.

For each question, the following information is recorded: question
number $q$, start time $t_{sq}$, end time $t_{eq}$, reaction time $t$,
and the selected answer option. These parameters are subsequently used
to synchronize with physiological signals and generate diagnostic
features.

The physiological data recording unit is designed for continuous
recording of an electrocardiogram during the psychological test and the
formation of a result file containing segmented ECG fragments for each
question.

The ECG signal is collected from the body surface using non-invasive
electrodes and fed to the analog-to-digital converter (ADC) of the
microprocessor module. The following steps are performed:

- preliminary analog filtering;

- analog-to-digital conversion;

- transmission of digital readings to the software subsystem.

The software registration module ensures continuous reception of digital
ECG readings and their storage in a text file of results in real time.

The results of electrocardiogram recording are saved in a text file of
ekg.txt format, in which the values of the ECG signal amplitude at
discrete moments in time are recorded line by line.

To ensure synchronization of physiological data with psychological
stimuli, a special \emph{separator marker is used}, inserted into the
file at the end of each question.

Each fragment between two markers corresponds to one test question.

The use of segmented ECG recordings divided by test questions provides
the following advantages:

- analysis of the body' s response to each stimulus;

- identification of hidden stress reactions;

- comparison of physiological parameters with specific psychological
scales;

- formation of an individual psychophysiological profile;

- increased diagnostic reliability.

This approach allows us to move from an integrated assessment of a
person's condition to event-oriented psychophysiological diagnostics.

One of the key features of the developed biotechnical
psychophysiological testing system is the use of event-driven
electrocardiogram recording during psychological testing tasks. Unlike
traditional physiological monitoring methods, which analyze the
subject' s holistic state throughout the entire
examination session, the proposed system utilizes event-driven
synchronization of physiological data with individual psychological
stimuli. Each test question is considered a distinct psychophysiological
event, eliciting an individual response from the body' s
regulatory systems. This approach allows for a shift from an average
assessment of functional status to a detailed analysis of responses to
specific cognitive and emotional stimuli, significantly increasing the
diagnostic sensitivity and prognostic value of the system.

Experimental testing of the developed biotechnical system for
psychophysiological testing was carried out with the aim of assessing
the performance of the hardware and software architecture, verifying the
algorithms for processing biomedical signals, and determining the
diagnostic information content of the proposed psychophysiological
indicators.

The main objective of the experimental studies was to confirm the
possibility of objectively recording the physiological reactions of the
body during the performance of psychological test tasks, as well as to
evaluate the effectiveness of the event-segmented approach to the
analysis of the electrocardiogram, synchronized with individual test
questions.

The experimental studies were conducted in a laboratory setting at the
Educational and Scientific Center for Psychophysiological Research.
Volunteers aged 18 to 30 with no diagnosed cardiovascular or
neurological conditions participated in the trial.

Before testing, subjects underwent a briefing explaining the
study' s objectives, the procedure, and the rules for
using the system' s hardware. Each participant provided
informed consent for the processing of personal and biomedical data.

The testing procedure included the following steps:

1. Registration of an individual background profile in a resting state
(3--5 minutes).

2. Conducting a battery of psychological tests with synchronous ECG
recording.

3. Formation of the resulting ECG file with segmentation by questions.

4. Automatic processing of physiological data.

5. Formation of a psychophysiological conclusion.

During the experiment, standard requirements of biomedical ergonomics
were observed: a comfortable position for the subject, a stable room
temperature, and the absence of external irritants and electromagnetic
interference.

During the testing, a battery of psychological tests was used, aimed at
diagnosing:

- anxiety level;

- aggressiveness of personality;

- neuropsychic stability;

- stress resistance and cognitive load.

Each test included 20 to 40 questions, presented in random order to
minimize habituation. For each question, the following were recorded:

- question number;

- the moment of beginning and end;

- selected answer option;

- the reaction time of the subject.

Simultaneously with the presentation of questions, continuous ECG
recording was carried out, followed by recording of data in a text file
of ekg.txt format with division of segments by marker.

The processing of experimental data was carried out automatically and
included the following steps:

1. Import of the resulting ECG file.

2. Signal segmentation by FFFF markers.

3. Digital filtering of each segment.

4. R-peak detection and cardiac cycle segmentation.

5. Calculation of RR intervals and heart rate variability indices.

6. Formation of a vector of psychophysiological characteristics.

7. Normalization relative to the background profile.

8. Intelligent classification of states.

9. Calculation of the integral index of psychophysiological stress.

For each test question, a set of indicators was formed that reflect the
body' s physiological response to the corresponding
stimulus.

The experiment yielded representative arrays of event-segmented
physiological data. Analysis revealed consistent patterns of changes in
cardiovascular regulation indicators depending on the emotional
significance of the stimuli.

For questions related to the anxiety and aggression scales, the
following were observed:

- increase in heart rate by 10-18\%;

- reduction of SDNN and RMSSD values by 15-25\%;

- growth of the integral energy of the cardiac signal;

- increasing the reaction time of the subject.

For neutral questions, physiological parameters remained close to
background values. The application of the ensemble classification model
provided the following average accuracy rates:

- Accuracy: 91- 93\%

- Sensitivity: 88 - 92\%

- Specificity: 94 - 97\%

- F1 score: 0.88-0.91

The obtained values confirm the high stability of the algorithms and the
correctness of the selected set of diagnostic features.

The results of the pilot study showed that the developed biotechnical
system can detect both overt and covert psycho-emotional reactions,
which are not always consciously recognized by the subjects. This
confirms the diagnostic value of the instrumental approach and its
advantages over traditional survey methods.

Event-segmented ECG analysis provides a detailed picture of the
body' s reactions to individual psychological stimuli and
allows for the formation of individual psychophysiological profiles
suitable for practical use in professional selection and functional
monitoring systems.

{\bfseries Conclusions.} This paper examines the design principles,
architecture, and practical application of a biotechnical system for
human psychophysiological testing. This system is designed to
objectively assess the functional and psychoemotional state of an
individual while completing psychological testing tasks. The developed
system implements an integrated biotechnical approach, combining an
automated psychological testing unit and a hardware unit for recording
physiological signals, followed by intelligent data processing.

The relevance of the research conducted is determined by the growing
role of the human factor in ensuring the reliability and safety of
modern technical and organizational systems {[}11-14{]}. Given the
increasing information load, the increasing complexity of professional
activities, and the increasing responsibility of personnel, traditional
psychodiagnostic methods, based primarily on survey methods, do not
always provide the necessary level of objectivity and predictive value.
Therefore, the development of psychophysiological testing tools based on
the recording of the body' s physiological reactions is
an important and sought-after scientific and technical task.

A key feature of the developed biotechnical system is the event-driven
organization of the testing procedure, in which each question in the
psychological test is treated as an individual psychophysiological
stimulus, eliciting a specific response from the body' s
regulatory systems. The psychological testing unit ensures standardized
presentation of stimulus material, recording of the
subject' s responses and the timing of task completion,
as well as the generation of a testing protocol containing data on the
sequence of questions, reaction times, and selected answers.

This study developed and tested a set of algorithms for digital
processing of biomedical signals, including multi-stage filtering,
artifact suppression, cardiac cycle segmentation, detection of
characteristic ECG points, heart rate variability analysis, and the
generation of an extended vector of diagnostic features. The use of an
event-based segmentation approach allows for the calculation of
psychophysiological indicators separately for each test question and the
comparison of physiological responses with specific psychological
scales.

The study places particular emphasis on individualizing diagnostics
through the use of a background profile of the subject, formed at rest.
Normalizing diagnostic features relative to the individual background
allows for accounting for interindividual variability in physiological
parameters and ensures accurate comparison of the same
individual' s responses under different conditions. This
approach improves the robustness of state recognition algorithms and the
reliability of diagnostic conclusions.

Experimental testing of the system confirmed the operability of the
hardware and software architecture, the correctness of the implemented
signal processing algorithms, and the high diagnostic value of the
event-based segmentation approach. The results demonstrate the ability
to detect both overt and covert psychoemotional reactions that are not
always consciously recognized by subjects and not reflected in their
subjective self-assessments. This confirms the fundamental advantages of
the instrumental psychophysiological approach compared to traditional
survey methods.

The practical significance of the developed biotechnical system lies in
its potential for use in professional selection systems and periodic
certification of personnel in technically complex and socially
responsible fields, such as aviation, transportation, energy, industry,
telecommunications, law enforcement, and the military. The mobility of
the hardware and the versatility of the software allow the system to be
used in both stationary laboratories and mobile settings, including
pre-shift and pre-trip inspections.

The use of event-segmented ECG recording during psychological testing
paves the way for the development of individual digital
psychophysiological profiles of employees and the implementation of
long-term monitoring systems for their functional status. Integration of
the system with corporate and departmental information platforms enables
the automation of data analysis, report generation, and decision support
processes in the areas of human resources management and occupational
safety.

Further development of the biotechnical system includes expanding the
sensor suite by incorporating galvanic skin response, respiratory
parameters, and accelerometry, which will improve diagnostic robustness
and the accuracy of psychophysiological response interpretation. The
implementation of deep learning methods for automatically generating
diagnostic features and adapting classifiers to the individual
characteristics of test subjects is also promising.

Thus, the developed biotechnical system for psychophysiological testing
represents a modern and effective tool for objectively diagnosing a
person' s functional and psychoemotional state.
Event-driven integration of psychological testing and segmented ECG
recording provides a fundamentally new level of detail and reliability
in psychophysiological diagnostics and forms the basis for the creation
of intelligent systems to ensure professional safety in the face of
increasing psychoemotional and informational stress.
\end{multicols}

\begin{center}
{\bfseries References}
\end{center}

\begin{refs}
1. Chernorizov AM, Isaychev SA, Zinchenko Yu. P. Psychophysiological
methods for the diagnostics of human functional states: New approaches
and perspectives // Psychology in Russia: State of the Art. -- 2016. -
Vol.9(4).- P.23-36. DOI 10.11621/ pir.2016.0403.

2. Lohani M, Payne BR, Strayer DL. A Review of Psychophysiological
Measures to Assess Cognitive States in Real-World Driving//Front Hum
Neurosci. - 2019. doi 10.3389/fnhum.2019.00057.

3. Tiberio L., Cesta A., Olivetti Belardinelli M. Psychophysiological
Methods to Evaluate User' s Response in Human Robot
Interaction: A Review and Feasibility Study. //Robotics.-
2013. -Vol.2(2). -P.92-121.
\href{https://doi.org/10.3390/robotics2020092}{DOI
10.3390/robotics2020092}.

4. Schmälzle R., Grall C. Psychophysiological Methods: Options, Uses,
and Validity. // The International Encyclopedia of Media Psychology, J.
Bulck (Ed.). - 2020.
\href{https://doi.org/10.1002/9781119011071.iemp0013}{DOI
10.1002/9781119011071.iemp0013}.

5. Li K, Rüdiger H, Ziemssen T. Spectral Analysis of Heart Rate
Variability: Time Window Matters// Front Neurol.- 2019.-Vol.10. doi
10.3389/ fneur.2019.00545.

6. Estévez M., Machado C., Leisman G., Estévez-Hernández T.,
Arias-Morales A., Machado A., Montes-Brown J. Spectral analysis of heart
rate variability //International Journal on Disability and Human
Development. - 2016. -Vol.15(1). - P.5-17.
\href{https://doi.org/10.1515/ijdhd-2014-0025}{DOI
10.1515/ijdhd-2014-0025}.

7. Pradhan BK, Neelappu B. Ch., Sivaraman J., Kim D., Pal K. A Review on
the Applications of Time-Frequency Methods in ECG Analysis// Journal of
Healthcare Engineering. - 2023.-Vol.1:3145483.
\href{https://doi.org/10.1155/2023/3145483}{DOI 10.1155/2023/3145483.}

8. Zyout A, Alquran H, Mustafa WA, Alqudah AM. Advanced Time-Frequency
Methods for ECG Waves Recognition. Diagnostics (Basel).
-2023.-Vol.1:308. doi 10.3390/diagnostics13020308.

9. Kei H., Satoshi Y., Yasuya I., Ryota Y., Ryo W., Naoki T., Noriyuki
S. Time-frequency analysis of high-frequency power profile within the
QRS complex in patients with ventricular tachycardia using wavelet
transform //Journal of Electrocardiology. -2025.-Vol.91:154027.
DOI 10.1016/j.jelectrocard.2025.154027.

10. Amekran Y., Damoun N., El H. Analysis of frequency-domain heart rate
variability using absolute versus normalized values: implications and
practical concerns // Frontiers in Physiology. -2024.-Vol.15:1470684.
DOI 10.3389/fphys.2024.1470684.

11. Mazakova A.T., Daribayeva G.D., Amirkhanov B.S., Zholmagambetova B.,
Abdirazak B.K. Biotechnical system of psychophysiological diagnostics of
personality conflict// Bulletin of KazUTB. - 2019. - No.4. - P.2-19.

12. Mazakova A.T., Amirkhanov B.S., Daribaeva G.D., Ziyatbekova G.Z.,
Abdirazak B.K. Программно-аппаратный комплекс психофизиологического
тестирования// Вестник КазУТБ. -2019.-№ 1.-С-2-9

13. Mazakova A. T., Amirkhanov B.S., Mussabek G.K.,
Abdirazak B.K., Daribaeva G.D. Expansion of psychological test battery
biotechnical psychophysiological diagnostic system // Bulletin of the
Kazakh-British Technical University. -2021.-Vol.18(3).-P.51-55. DOI
10.55452/1998-6688-2021-18-3-51-55.

14. A.s. No.5499, 2019, September 26. Mazakov T.Z., Ziyatbekova G.Z.,
Mazakova A.T., Daribaeva G.D., Əbdirazak B.K., Amirkhanov B.S.,
Zholmagambetova B.R. A set of psychophysiological testing programs.
Computer program.
\end{refs}

\begin{info}
\hspace{1em}\emph{{\bfseries Information about the authors}}

Mazakova A. T. - PhD, senior lecturer, Al-Farabi Kazakh National
University, Almaty, Kazakhstan, e-mail:
aigerym97@mail.ru;

Jomartova Sh.А.- Doctor of Technical Sciences, Professor, Al-Farabi
Kazakh National University, Almaty, Kazakhstan, e-mail:
jomartova@mail.ru;

Mazakov T. Zh. - Doctor of Physical and Mathematical Sciences,
Professor, Al-Farabi Kazakh National University and International
University of Engineering and Technology, Almaty, Kazakhstan, e-mail:
tmazakov@mail.ru;

Zhao Baidong - Senior Lecturer, Dalian Polytechnic University,
Dalian, China, e-mail: xiaoyuezhishang@vip.qq.com;

Abubakirov A. M. - 2nd year Master, Al-Farabi Kazakh National
University, Almaty, Kazakhstan, e-mail: adlett28@gmail.com.

\hspace{1em}\emph{{\bfseries Сведения об авторах}}

Мазакова А.N. - PhD, старший преподаватель КазНУ им. аль-Фараби, Алматы,
Казахстан, e-mail:
aigerym97@mail.ru;

Джомартова Ш. А. -- д.т.н., профессор КазНУ им. аль-Фараби, Алматы,
Казахстан, e-mail:
jomartova@mail.ru;

Мазаков Т. Ж. - д.ф.-м.н., КазНУ им. аль-Фараби, Международный
инженерно-технологический университет, Алматы, Казахстан, e-mail:
tmazakov@mail.ru;

Чжао Байдун - старший преподаватель, Даляньский политехнический
университет, Далянь, Китай, e-mail: xiaoyuezhishang@vip.qq.com;

Абубакиров А. М. - 2 курс магистратуры, КазНУ им. Аль-Фараби, Алматы,
Казахстан, e-mail:
adlett28@gmail.com.
\end{info}
